% !TeX root = ../papers/nc-p3.tex
\section{非対称成分の反対行列式と二元障害を切る残余因子}
\label{sec:binary-opposite-determinant}
\begin{theorem}\label{thm:binary-opposite-determinant}
前定理の15次元成分の、同じ有理標本と二元障害の一般点での軌道横断局所環を取る。
二つの片側対称因子の式を $a,b$ とすると、反対切断の行列式は単元倍を除き
\[
                         a^6b^6 D
\]
となる。$a,b,D$ は正則パラメータ系の一部で、$a=b=0$ において
$D=0$ は指定の向きの二元障害の閉包と局所的に一致する。
15次元成分の一般反対階数は16、二つの片側対称因子の一般反対階数はともに10、
残余因子の一般反対階数は15である。残余因子の一般点は有限被約グラフを持ち、
両方の反対交換対称性を欠く最小周期4の安定なFrobenius窓である。
この階数15の因子自身の一般点のAS正則性は、周囲の一般点の正則性からは決まらない。
\end{theorem}
\begin{proof}[因子の抽出と二元錐の一次関係]
局所模型の混合座標をそれぞれ指標 $+-$ と $-+$ の $a,b$ とし、
テンソルの各反対組を対称十成分と交代六成分に分ける。
同変な陰関数消去で戻した反対行列は
\[
           \begin{pmatrix} O&aP\\ bQ&abF\end{pmatrix}
\]
で、四ブロック $O,P,Q,F$ は二交換で不変である。交代六列から $a$、交代六行から $b$ をくくると
残余行列 $M=\left(\begin{smallmatrix}O&P\\Q&F\end{smallmatrix}\right)$ とその行列式 $D$ が出る。

指定の側が $U=P_2\oplus kz$ の二元点では、錐の二次式
$C=\Delta(x,x')^2\in\operatorname{Sym}^2 P_2$ が $O$ の左核に入る。
この式は $P$ の値も消す。実際 $P$ は $+-$ の接変分の対称・交代ブロックであり、
接条件からその変分は $K_0(x,y)\otimes K_0(x',y')$ と片側を交換した空間の双方を消す。
ここで $K_0=\Delta(x,y)(P_1(x)\otimes P_2(y))$ である。
$\bigwedge^2 P_3(y)$ の元は $\Delta(y,y')$ と対称な双次数 $(2,2)$ の式の積である。
二次元のPlücker恒等式
\[
 \Delta(x,x')\Delta(y,y')
 =\Delta(x,y)\Delta(x',y')-\Delta(x,y')\Delta(x',y)
\]
に残る $\Delta(x,x')$ と双次数 $(2,2)$ の式を掛ければ、
$C\otimes\bigwedge^2P_3(y)$ は上記の消える二空間の和に入る。
従って $M$ の左核に $(C,0)$ があり、二元障害の上で $D=0$ となる。
\end{proof}
\begin{proof}[交叉に沿った残余行列式の微分]
計算する成分と持上げを明示する。$A,B$ を二つの混合接の非軌道代表とし、
$X$ を $a=b=0$ の13次元交叉に接する対称方向とする。第一切断の一般逆を $J$、
核列を $N$ とし、$[V,W]=VJW+WJV$ と略記する。
$M$ の定数項では $O,P,Q$ はそれぞれ $w,A,B$ の反対切断の該当ブロック、
$F$ は混合項 $Z$ の交代・交代ブロックである。$Z$ は $--$ の空間で
$N^t ZN=N^t[A,B]N$ の唯一の解である。

次に方向 $X$ でこれらを微分する。$A_1,B_1$ は対応する混合指標で
$N^t A_1 N=N^t[X,A]N$, $N^t B_1N=N^t[X,B]N$ を解く。
$Z_1$ は $--$ の空間で、両辺に $N^t, N$ を掛けたとき次の式の右辺に一致する解である。
\[
 Z_1\ \equiv\ [X,Z]+[A,B_1]+[B,A_1]
          -\sum_{(U,V,W)\in\operatorname{Perm}(X,A,B)} UJVJW .
\]
これはSchur補の三次までの展開で、$X,A,B$ の三方向にそれぞれ一次の係数を取った式である。
$M$ の微分の四ブロックは $X,A_1,B_1,Z_1$ から取れる。
解の存在は前定理の滑らかな15次元枝で $X$ を交叉に接して取ったことによる。
混合の一次補正の解は自由な混合接を加えられるが、これは混合座標の一階の規格化と
文字基底の一階変更で実現できるので任意の一組でよい。

\path{computations/verify_binary_opposite_determinant.py}が有理標本で素数1009により次を検査する。
定数 $O$ は階数8、その左右二次元核に制限した $P,Q$ の階数は順に1と2、
$M$ の階数は15、ある $X$ に対する微分を $M$ の左右核線で縮約した値は非零である。
上の錐の恒等式が $M$ の階数の上界15を与える。
線型系の全Jacobi階数は前定理の有理数検証と一致し、
$X$ に課す二つの混合障害は標数零でも同じ非零の微分を与える。
従って各接と補正をこの良い素数での指定解に還元するよう持ち上げられ、非零微分は標数零でも成り立つ。
これで $D$ が $a=b=0$ 上でも正則パラメータとなる。
その零点は滑らかな12次元で、そこに入る既約な12次元の二元障害の閉包は局所的に全体と一致する。
微分の非消滅と階数の小行列は開条件なので、二元障害の一般点にも同じ結論が延びる。
\end{proof}
\begin{proof}[反対階数]
$a,b,D$ は独立だから三因子以外では反対階数16である。
残余因子の一般点は $ab\ne0$ で $M$ は階数15のままなので、実際の反対階数も15。
被約有限グラフと最小周期については、一般二元点で前定理の番号付けを取り、
その点を通る $ab\ne0$ の部分に同じ置換の特殊化を適用すればよい。
有限点グラフでの安定性も一般二元点から開条件で従う。

片側対称因子上では階数は10以下で、上記の $P,Q$ の非零制限からそれぞれ9以上、10以上である。
中央支持3以下の点の存在は射影的な入射から閉条件で、有理標本は非正則なのでこの条件を満たさない
（定理\ref{thm:finite-cubic-separators}）。従って両因子の一般点でもこの条件はない。
ところが定理\ref{thm:finite-rank-nine-all-marks}により、有限グラフで階数9の正則窓には
中央支持3の点がある。前定理の両因子の一般正則性から階数9が除かれ、いずれも10となる。
\end{proof}
